\documentclass[12pt,a4paper]{article}

% ==================== PACKAGES ====================
\usepackage[utf8]{inputenc}
\usepackage[vietnamese]{babel}
\usepackage{amsmath,amssymb,amsthm}
\usepackage{geometry}
\usepackage{enumitem}
\usepackage[most]{tcolorbox}
\usepackage{hyperref}
\usepackage{fancyhdr}
\usepackage{graphicx}
\usepackage{tikz}
\usepackage{eso-pic}
\usepackage{xcolor}

\geometry{a4paper, margin=2cm, bottom=2.5cm}
\hypersetup{colorlinks=true, linkcolor=blue!70!black, urlcolor=blue}
\setlist[itemize]{topsep=4pt, itemsep=2pt}
\setlist[enumerate]{topsep=4pt, itemsep=2pt}

% ==================== WATERMARK ====================
\AddToShipoutPictureFG{%
  \AtPageCenter{%
    \begin{tikzpicture}[remember picture, overlay]
      \node[rotate=45, scale=1, opacity=0.06]{%
        \fontsize{60}{72}\selectfont\bfseries\sffamily Đặng Tấn Quí%
      };
    \end{tikzpicture}%
  }%
}

% ==================== HEADER / FOOTER ====================
\pagestyle{fancy}
\fancyhf{}
\fancyhead[L]{\small\textit{Toán cao cấp 1 -- Đại số tuyến tính}}
\fancyhead[R]{\small\textit{Đặng Tấn Quí}}
\fancyfoot[C]{\thepage}
\fancyfoot[L]{\footnotesize\textcopyright\ Đặng Tấn Quí}
\fancyfoot[R]{\footnotesize SĐT: 0377\,122\,210}
\renewcommand{\headrulewidth}{0.4pt}
\renewcommand{\footrulewidth}{0.4pt}

\fancypagestyle{plain}{%
  \fancyhf{}
  \fancyfoot[C]{\thepage}
  \fancyfoot[L]{\footnotesize\textcopyright\ Đặng Tấn Quí}
  \fancyfoot[R]{\footnotesize SĐT: 0377\,122\,210}
  \renewcommand{\headrulewidth}{0pt}
  \renewcommand{\footrulewidth}{0.4pt}
}

% ==================== CUSTOM ENVIRONMENTS ====================
\newtcolorbox{dinhnghia}[1][]{%
  enhanced, breakable,
  colback=blue!3!white, colframe=blue!60!black,
  fonttitle=\bfseries, title={Định nghĩa}, #1%
}

\newtcolorbox{dinhly}[1][]{%
  enhanced, breakable,
  colback=green!3!white, colframe=green!50!black,
  fonttitle=\bfseries, title={Định lý}, #1%
}

\newtcolorbox{tinhchat}[1][]{%
  enhanced, breakable,
  colback=yellow!5!white, colframe=orange!50!black,
  fonttitle=\bfseries, title={Tính chất}, #1%
}

\newtcolorbox{phuongphap}[1][]{%
  enhanced, breakable,
  colback=gray!5!white, colframe=gray!50!black,
  fonttitle=\bfseries, title={Phương pháp}, #1%
}

\newtcolorbox{luuy}[1][]{%
  enhanced, breakable,
  colback=red!3!white, colframe=red!50!black,
  fonttitle=\bfseries, title={Lưu ý}, #1%
}

\newtcolorbox{congthuc}[1][]{%
  enhanced, breakable,
  colback=cyan!3!white, colframe=cyan!50!black,
  fonttitle=\bfseries, title={Công thức}, #1%
}

\newtcolorbox{vidu}[1][]{%
  enhanced, breakable,
  colback=violet!3!white, colframe=violet!50!black,
  fonttitle=\bfseries, title={Ví dụ}, #1%
}

% ==================== DOCUMENT ====================
\begin{document}

% ==================== TRANG BÌA ====================
\thispagestyle{empty}
\begin{tikzpicture}[remember picture, overlay]
  \fill[blue!80!black] (current page.south west) rectangle (current page.north east);
  \fill[blue!60!cyan, opacity=0.8]
    (current page.south west) -- (current page.north west) --
    ([xshift=-3cm]current page.north east) -- (current page.south east) -- cycle;
  \foreach \r/\o in {6/0.05, 4.5/0.08, 3/0.12} {
    \fill[white, opacity=\o] ([xshift=2cm, yshift=-2cm]current page.north east) circle (\r cm);
  }
  \draw[white, line width=2pt] ([yshift=-5cm]current page.north west) ++(2cm,0) -- ++(17cm,0);
  \draw[white, line width=2pt] ([yshift=-16cm]current page.north west) ++(2cm,0) -- ++(17cm,0);
  \node[anchor=west, white, font=\fontsize{28}{34}\bfseries\selectfont]
    at ([xshift=3cm, yshift=-7.5cm]current page.north west)
    {TOÁN CAO CẤP 1};
  \node[anchor=west, white, font=\fontsize{20}{26}\selectfont]
    at ([xshift=3cm, yshift=-9cm]current page.north west)
    {Đại số tuyến tính --- Năm nhất đại học};
  \node[anchor=west, white, font=\fontsize{18}{24}\selectfont]
    at ([xshift=3cm, yshift=-10cm]current page.north west)
    {Tài liệu lý thuyết};
  \node[anchor=west, white, font=\Large\bfseries]
    at ([xshift=3cm, yshift=-13cm]current page.north west) {Biên soạn:};
  \node[anchor=west, yellow!80!white, font=\fontsize{24}{30}\bfseries\selectfont]
    at ([xshift=3cm, yshift=-14.5cm]current page.north west) {ĐẶNG TẤN QUÍ};
  \node[anchor=west, white!90, font=\large]
    at ([xshift=3cm, yshift=-18cm]current page.north west) {SĐT: 0377\,122\,210};
  \node[anchor=south, white!80, font=\small]
    at ([yshift=1.5cm]current page.south) {Tài liệu có bản quyền --- Cấm sao chép khi chưa được sự đồng ý của tác giả};
  \node[anchor=south east, white, font=\Large\bfseries]
    at ([xshift=-2cm, yshift=3cm]current page.south east) {\the\year};
\end{tikzpicture}

\newpage
\tableofcontents
\newpage

%% ================================================================
%%  CHƯƠNG 1: MA TRẬN VÀ ĐỊNH THỨC
%% ================================================================
\section{Ma trận}

\subsection{Định nghĩa ma trận}

\begin{dinhnghia}
  \textbf{Ma trận} cỡ $m \times n$ (hoặc cấp $m \times n$) là bảng số hình chữ nhật gồm $m$ dòng và $n$ cột:
  \[
    A = \begin{pmatrix}
      a_{11} & a_{12} & \cdots & a_{1n} \\
      a_{21} & a_{22} & \cdots & a_{2n} \\
      \vdots & \vdots & \ddots & \vdots \\
      a_{m1} & a_{m2} & \cdots & a_{mn}
    \end{pmatrix} = (a_{ij})_{m \times n}.
  \]
  Phần tử $a_{ij}$ nằm ở dòng $i$, cột $j$.
\end{dinhnghia}

\begin{dinhnghia}
  \textbf{Ma trận vuông} cấp $n$: $m = n$. Đường chéo chính: $a_{11}, a_{22}, \ldots, a_{nn}$.
\end{dinhnghia}

\begin{tinhchat}[title={Các loại ma trận đặc biệt}]
  \begin{itemize}
    \item \textbf{Ma trận không} $\mathbf{0}$: mọi phần tử bằng 0.
    \item \textbf{Ma trận chuyển vị} $A^T$: $(A^T)_{ij} = a_{ji}$. Đổi hàng và cột cho nhau.
    \item \textbf{Ma trận đối xứng}: $A^T = A$ (ma trận vuông).
    \item \textbf{Ma trận tam giác trên}: $a_{ij} = 0$ khi $i > j$ (phần tử dưới đường chéo bằng 0).
    \item \textbf{Ma trận tam giác dưới}: $a_{ij} = 0$ khi $i < j$ (phần tử trên đường chéo bằng 0).
    \item \textbf{Ma trận chéo}: $a_{ij} = 0$ khi $i \ne j$ (chỉ có đường chéo chính khác 0).
    \item \textbf{Ma trận bằng nhau}: $A = B$ khi và chỉ khi cùng cỡ và $a_{ij} = b_{ij}$ với mọi $i, j$.
  \end{itemize}
\end{tinhchat}

\subsection{Các phép toán trên ma trận}

\begin{dinhnghia}
  \textbf{Cộng ma trận} (cùng cỡ): $(A + B)_{ij} = a_{ij} + b_{ij}$.
\end{dinhnghia}

\begin{tinhchat}
  \begin{itemize}
    \item $A + B = B + A$
    \item $A + 0 = A$
    \item $A + (-A) = 0$
    \item $(A + B) + C = A + (B + C)$
  \end{itemize}
\end{tinhchat}

\begin{dinhnghia}
  \textbf{Nhân ma trận với số}: $(\lambda A)_{ij} = \lambda a_{ij}$.
\end{dinhnghia}

\begin{tinhchat}
  \begin{itemize}
    \item $k(A + B) = kA + kB$
    \item $(h + k)A = hA + kA$
    \item $1.A = A$
    \item $0.A = 0$
  \end{itemize}
\end{tinhchat}

\begin{dinhnghia}
  \textbf{Nhân ma trận} $A_{m \times p}$ với $B_{p \times n}$ ta được ma trận $C_{m \times n}$:
  \[
    c_{ij} = (AB)_{ij} = \sum_{k=1}^{p} a_{ik} b_{kj}.
  \]
  Điều kiện: số cột của $A$ = số dòng của $B$.
\end{dinhnghia}

\begin{tinhchat}
  \begin{itemize}
    \item $A(BC) = (AB)C$.
    \item $A(B + C) = AB + AC$.
    \item $(B + C)A = BA + CA$.
    \item $k(BC) = (kB)C = B(kC)$
    \item $(AB)^T = B^T A^T$.
    \item $AB$ được chưa chắc $BA$ được. Nếu ma trận vuông thì $AB$ và $BA$ được
    \item Nói chung $AB \ne BA$.
    \item $A \ne 0, B \ne 0$ nhưng có thể $AB = 0$
  \end{itemize}
\end{tinhchat}

\subsection{Định thức}

\begin{dinhnghia}
  \textbf{Định thức} của ma trận vuông $A$ cấp $n$, ký hiệu $\det(A)$ hoặc $|A|$:
  \begin{itemize}
    \item Cấp 1: $\det(a) = a$.
    \item Cấp 2: $\det\begin{pmatrix} a & b \\ c & d \end{pmatrix} = ad - bc$.
    \item Cấp $n$: khai triển theo dòng $i$: $\det(A) = \sum_{j=1}^{n} (-1)^{i+j} a_{ij} M_{ij}$, với $M_{ij}$ là định thức con bù của $a_{ij}$ là bỏ hàng $i$ và cột $j$.
  \end{itemize}
\end{dinhnghia}

\begin{dinhly}[title={Tính chất định thức}]
  \begin{itemize}
    \item $\det(A^T) = \det(A)$.
    \item Đổi chỗ hai dòng (hoặc hai cột) $\Rightarrow$ định thức đổi dấu.
    \item Nhân một dòng (cột) với $\lambda$ $\Rightarrow$ định thức nhân với $\lambda$.
    \item $\det(AB) = \det(A) \cdot \det(B)$.
    \item Ma trận có 2 hàng (cột) giống nhau $\Rightarrow$ $\det(A) = 0$.
    \item Ma trận có 2 hàng (cột) tỉ lệ $\Rightarrow$ $\det(A) = 0$.
    \item Ma trận có 1 hàng (cột) bằng 0 $\Rightarrow$ $\det(A) = 0$.
    \item Ma trận có 1 hàng (cột) là tổ hợp tuyến tính của các hàng khác hoặc các cột khác $\Rightarrow$ $\det(A) = 0$.
    \item Ma trận có tất cả các phần tử của 1 hàng (cột) có dạng tổng 2 số hạng thì có thể tách thành tổng 2 định thức.
    \item Khi cộng bội k của 1 hàng vào 1 hàng khác (hay bội k của 1 cột vào 1 cột khác) thì được định thức mới bằng định thức cũ.
    \item Ma trận tam giác có định thức là tích đường chéo
  \end{itemize}
\end{dinhly}

\begin{phuongphap}[title={Tính định thức}]
  \begin{enumerate}
    \item Khai triển theo dòng hoặc cột có nhiều số 0.
    \item Biến đổi sơ cấp đưa về ma trận tam giác (định thức = tích phần tử đường chéo).
    \begin{enumerate}
      \item Nhân 1 hàng với 1 số $k \ne 0$
      \item Đổi chỗ 2 hàng
      \item Cộng $k$ lần hàng này vào hàng kia
    \end{enumerate}
  \end{enumerate}
\end{phuongphap}

\subsection{Ma trận nghịch đảo}

\begin{dinhnghia}
  \begin{itemize}
    \item \textbf{Ma trận đơn vị} $I_n$: ma trận vuông cấp $n$, đường chéo chính bằng 1, còn lại bằng 0. Ký hiệu $(I_n)_{ij} = \delta_{ij}$ (Kronecker delta).
    \item $I \cdot A = A \cdot I = A$ với mọi ma trận vuông $A$ cùng cấp.
  \end{itemize}
\end{dinhnghia}

\begin{dinhnghia}
  Ma trận vuông $A$ gọi là \textbf{khả đảo} (khả nghịch) nếu tồn tại ma trận $B$ sao cho $AB = BA = I$. Khi đó $B$ gọi là \textbf{ma trận nghịch đảo} của $A$, ký hiệu $A^{-1}$. Ma trận nghịch đảo nếu tồn tại thì duy nhất.
\end{dinhnghia}

\begin{dinhly}[title={Điều kiện khả đảo}]
  Ma trận vuông $A$ khả đảo khi và chỉ khi $\det(A) \ne 0$.
\end{dinhly}

\begin{congthuc}[title={Công thức nghịch đảo qua phụ đại số}]
  Gọi $A_{ij}$ là phần bù đại số của $a_{ij}$ (tức $A_{ij} = (-1)^{i+j} M_{ij}$ với $M_{ij}$ là định thức con bù). Ma trận phụ hợp $\mathrm{adj}(A)$ có phần tử $(i,j)$ là $A_{ji}$. Khi đó:
  \[
    A^{-1} = \frac{1}{\det(A)} \cdot \mathrm{adj}(A).
  \]
  Với ma trận cấp 2: $A = \begin{pmatrix} a & b \\ c & d \end{pmatrix}$, $\det(A) \ne 0$ thì
  $A^{-1} = \dfrac{1}{ad-bc} \begin{pmatrix} d & -b \\ -c & a \end{pmatrix}$.
\end{congthuc}

\begin{tinhchat}[title={Tính chất ma trận nghịch đảo}]
  \begin{itemize}
    \item $(A^{-1})^{-1} = A$.
    \item $(AB)^{-1} = B^{-1} A^{-1}$ (nghịch đảo của tích = tích nghịch đảo theo thứ tự ngược).
    \item $(A^T)^{-1} = (A^{-1})^T$.
    \item $\det(A^{-1}) = \dfrac{1}{\det(A)}$.
  \end{itemize}
\end{tinhchat}

\begin{phuongphap}[title={Tính $A^{-1}$ bằng Gauss-Jordan}]
  Viết ma trận mở rộng $[A \mid I]$. Dùng biến đổi sơ cấp trên hàng đưa $A$ về $I$. Khi đó $I$ biến thành $A^{-1}$: $[A \mid I] \xrightarrow{\text{BĐSC}} [I \mid A^{-1}]$.
\end{phuongphap}

\subsection{Hạng của ma trận}

\begin{dinhnghia}
  \textbf{Hạng} của ma trận $A$, ký hiệu $\mathrm{rank}(A)$ hoặc $r(A)$, là cấp cao nhất của định thức con khác 0 của $A$. Nói cách khác: $\mathrm{rank}(A)$ = số dòng (hoặc cột) độc lập tuyến tính tối đa.
\end{dinhnghia}

\begin{tinhchat}[title={Tính chất hạng}]
  \begin{itemize}
    \item $\mathrm{rank}(A) = \mathrm{rank}(A^T)$.
    \item $\mathrm{rank}(A) \le \min(m, n)$ với $A$ cỡ $m \times n$.
    \item $\mathrm{rank}(AB) \le \min(\mathrm{rank}(A), \mathrm{rank}(B))$.
    \item $A$ khả đảo $\Leftrightarrow$ $\mathrm{rank}(A) = n$ (với $A$ vuông cấp $n$).
  \end{itemize}
\end{tinhchat}

\begin{phuongphap}[title={Tính hạng bằng biến đổi sơ cấp}]
  Dùng biến đổi sơ cấp trên hàng (hoặc cột) đưa $A$ về dạng bậc thang. Số dòng khác 0 của ma trận bậc thang chính là $\mathrm{rank}(A)$. Các biến đổi sơ cấp không làm thay đổi hạng.
\end{phuongphap}

%% ================================================================
%%  CHƯƠNG 2: HỆ PHƯƠNG TRÌNH TUYẾN TÍNH
%% ================================================================
\section{Hệ phương trình tuyến tính}

\subsection{Dạng ma trận}

\begin{dinhnghia}
  \textbf{Hệ phương trình tuyến tính} $m$ phương trình $n$ ẩn:
  \[
    \begin{cases}
      a_{11}x_1 + a_{12}x_2 + \cdots + a_{1n}x_n = b_1 \\
      a_{21}x_1 + a_{22}x_2 + \cdots + a_{2n}x_n = b_2 \\
      \vdots \\
      a_{m1}x_1 + a_{m2}x_2 + \cdots + a_{mn}x_n = b_m
    \end{cases}
  \]
  Viết dạng ma trận: $Ax = b$, trong đó
  \[
    A = \begin{pmatrix} a_{11} & \cdots & a_{1n} \\ \vdots & \ddots & \vdots \\ a_{m1} & \cdots & a_{mn} \end{pmatrix},\quad
    x = \begin{pmatrix} x_1 \\ \vdots \\ x_n \end{pmatrix},\quad
    b = \begin{pmatrix} b_1 \\ \vdots \\ b_m \end{pmatrix}.
  \]
  $A$: \textbf{ma trận hệ số} (cỡ $m \times n$). $x$: \textbf{vectơ ẩn}. $b$: \textbf{vectơ vế phải}. \textbf{Ma trận mở rộng} $\overline{A} = [A \mid b]$ (thêm cột $b$ vào bên phải $A$).
\end{dinhnghia}

\begin{dinhnghia}
  Hệ gọi là \textbf{thuần nhất} nếu $b = 0$ (tức $b_1 = \cdots = b_m = 0$). Hệ thuần nhất luôn có ít nhất nghiệm tầm thường $x = 0$.
\end{dinhnghia}

\subsection{Phương pháp Gauss và Gauss-Jordan}

\begin{phuongphap}[title={Phương pháp Gauss}]
  Dùng \textbf{biến đổi sơ cấp trên hàng} đưa ma trận mở rộng $\overline{A}$ về dạng \textbf{bậc thang} (tam giác trên):
  \begin{enumerate}
    \item Chọn phần tử khác 0 ở cột đầu (pivot), đưa lên hàng đầu nếu cần.
    \item Dùng pivot khử các phần tử bên dưới trong cùng cột (cộng bội của hàng chứa pivot).
    \item Lặp với ma trận con (bỏ hàng và cột đã xử lý).
  \end{enumerate}
  Từ ma trận bậc thang, thay thế ngược từ dòng cuối lên để tìm nghiệm.
\end{phuongphap}

\begin{phuongphap}[title={Phương pháp Gauss-Jordan}]
  Sau khi đưa về dạng bậc thang, tiếp tục biến đổi để đưa phần khác 0 của ma trận hệ số về \textbf{dạng ma trận đơn vị rút gọn} (mỗi cột pivot chỉ có 1 số 1, còn lại là 0). Khi đó nghiệm đọc trực tiếp từ cột vế phải.
  
  \textbf{Ứng dụng tính $A^{-1}$}: Viết $[A \mid I]$, dùng Gauss-Jordan đưa $A$ về $I$. Khi đó $I$ biến thành $A^{-1}$.
\end{phuongphap}

\begin{luuy}
  Ba phép biến đổi sơ cấp trên hàng: (1) Đổi chỗ hai hàng. (2) Nhân một hàng với số $k \ne 0$. (3) Cộng $k$ lần hàng này vào hàng kia. Các phép này không đổi tập nghiệm của hệ.
\end{luuy}

\subsection{Hệ Cramer}

\begin{dinhly}[title={Công thức Cramer}]
  Hệ $n$ phương trình $n$ ẩn $Ax = b$. Nếu $\det(A) \ne 0$ thì hệ có \textbf{nghiệm duy nhất}:
  \[
    x_j = \frac{\det(A_j)}{\det(A)}, \quad j = 1, \ldots, n,
  \]
  với $A_j$ là ma trận $A$ thay cột $j$ bởi cột $b$.
\end{dinhly}

\begin{vidu}
  Hệ $\begin{cases} 2x_1 + x_2 = 5 \\ x_1 + 3x_2 = 6 \end{cases}$. $\det(A) = 5$, $\det(A_1) = 9$, $\det(A_2) = 7$. Vậy $x_1 = 9/5$, $x_2 = 7/5$.
\end{vidu}

\subsection{Hệ thuần nhất}

\begin{dinhly}[title={Hệ thuần nhất $Ax = 0$}]
  \begin{itemize}
    \item Luôn có nghiệm $x = 0$ (nghiệm tầm thường).
    \item Có nghiệm không tầm thường khi và chỉ khi $\mathrm{rank}(A) < n$ (tức $A$ có $n$ cột phụ thuộc tuyến tính).
    \item Khi $\det(A) = 0$: hệ có vô số nghiệm; không gian nghiệm có số chiều $n - \mathrm{rank}(A)$.
  \end{itemize}
\end{dinhly}

\subsection{Định lý Kronecker-Capelli}

\begin{dinhly}[title={Định lý Kronecker-Capelli}]
  Hệ $Ax = b$ có nghiệm khi và chỉ khi $\mathrm{rank}(A) = \mathrm{rank}(\overline{A})$.
  \begin{itemize}
    \item Nếu $\mathrm{rank}(A) = \mathrm{rank}(\overline{A}) = n$: hệ có \textbf{nghiệm duy nhất}.
    \item Nếu $\mathrm{rank}(A) = \mathrm{rank}(\overline{A}) = r < n$: hệ có \textbf{vô số nghiệm}; số ẩn tự do $= n - r$; nghiệm tổng quát có dạng vectơ riêng + tổ hợp tuyến tính của $n-r$ vectơ cơ sở của không gian nghiệm hệ thuần nhất.
    \item Nếu $\mathrm{rank}(A) < \mathrm{rank}(\overline{A})$: hệ \textbf{vô nghiệm}.
  \end{itemize}
\end{dinhly}

%% ================================================================
%%  CHƯƠNG 3: KHÔNG GIAN VECTƠ
%% ================================================================
\section{Không gian vectơ}

\subsection{Định nghĩa và tiên đề}

\begin{dinhnghia}
  Tập $V$ với hai phép toán: \textbf{cộng} $+: V \times V \to V$ và \textbf{nhân vô hướng} $\cdot: \mathbb{R} \times V \to V$, gọi là \textbf{không gian vectơ} trên $\mathbb{R}$ nếu thỏa:
  \begin{enumerate}
    \item $(V, +)$ là nhóm Abel: giao hoán, kết hợp, có phần tử không $\mathbf{0}$, mỗi $v$ có phần tử đối $-v$.
    \item Phân phối: $\lambda(u + v) = \lambda u + \lambda v$, $(\lambda + \mu)v = \lambda v + \mu v$.
    \item Kết hợp vô hướng: $(\lambda\mu)v = \lambda(\mu v)$.
    \item Đơn vị: $1 \cdot v = v$.
  \end{enumerate}
\end{dinhnghia}

\begin{vidu}
  $\mathbb{R}^n$ (vectơ cột $n$ chiều); $M_{m \times n}(\mathbb{R})$ (ma trận $m \times n$); tập các đa thức bậc $\le n$ với hệ số thực; $C[a,b]$ (hàm liên tục trên $[a,b]$).
\end{vidu}

\subsection{Không gian con}

\begin{dinhnghia}
  Tập con $W \subset V$ gọi là \textbf{không gian con} của $V$ nếu $W$ đóng với phép cộng và nhân vô hướng: $u, v \in W$, $\lambda \in \mathbb{R}$ $\Rightarrow$ $u + v \in W$, $\lambda u \in W$. Khi đó $W$ cũng là KGVT với các phép toán cảm sinh.
\end{dinhnghia}

\begin{phuongphap}[title={Kiểm tra không gian con}]
  $W \ne \varnothing$ và $\forall u, v \in W$, $\forall \lambda \in \mathbb{R}$: $u + v \in W$, $\lambda u \in W$. Hoặc tương đương: $\forall u, v \in W$, $\forall \alpha, \beta \in \mathbb{R}$: $\alpha u + \beta v \in W$.
\end{phuongphap}

\subsection{Tổ hợp tuyến tính và tập sinh}

\begin{dinhnghia}
  Vectơ $v$ là \textbf{tổ hợp tuyến tính} của $\{v_1, \ldots, v_k\}$ nếu $v = \lambda_1 v_1 + \cdots + \lambda_k v_k$ với $\lambda_i \in \mathbb{R}$.
\end{dinhnghia}

\begin{dinhnghia}
  \textbf{Bao tuyến tính} (hay không gian con sinh bởi): $\mathrm{span}\{v_1, \ldots, v_k\} = \{\lambda_1 v_1 + \cdots + \lambda_k v_k \mid \lambda_i \in \mathbb{R}\}$. Hệ $\{v_1, \ldots, v_k\}$ \textbf{sinh} $V$ nếu $\mathrm{span}\{v_1, \ldots, v_k\} = V$.
\end{dinhnghia}

\subsection{Độc lập tuyến tính, phụ thuộc tuyến tính}

\begin{dinhnghia}
  Hệ vectơ $\{v_1, \ldots, v_k\}$ \textbf{phụ thuộc tuyến tính} nếu tồn tại $\lambda_1, \ldots, \lambda_k$ không đồng thời bằng 0 sao cho $\lambda_1 v_1 + \cdots + \lambda_k v_k = \mathbf{0}$. Ngược lại là \textbf{độc lập tuyến tính}.
\end{dinhnghia}

\begin{tinhchat}
  Hệ chứa vectơ $\mathbf{0}$ thì phụ thuộc tuyến tính. Một vectơ $v \ne \mathbf{0}$ thì độc lập tuyến tính. Hai vectơ phụ thuộc tuyến tính $\Leftrightarrow$ chúng tỉ lệ.
\end{tinhchat}

\begin{dinhnghia}
  \textbf{Hạng} của hệ vectơ $\mathrm{rank}\{v_1, \ldots, v_k\}$: số vectơ tối đa độc lập tuyến tính trong hệ (bằng số chiều của $\mathrm{span}\{v_1, \ldots, v_k\}$).
\end{dinhnghia}

\subsection{Cơ sở và số chiều}

\begin{dinhnghia}
  Hệ vectơ $\mathcal{B} = \{e_1, \ldots, e_n\}$ là \textbf{cơ sở} của $V$ nếu $\mathcal{B}$ độc lập tuyến tính và sinh $V$ (mọi $v \in V$ biểu diễn duy nhất: $v = x_1 e_1 + \cdots + x_n e_n$).
\end{dinhnghia}

\begin{dinhly}
  Mọi cơ sở của $V$ có cùng số phần tử. Số đó gọi là \textbf{số chiều} của $V$, ký hiệu $\dim V$. Nếu $\dim V = n$ thì mọi hệ $n$ vectơ độc lập tuyến tính đều là cơ sở; mọi hệ sinh có ít nhất $n$ phần tử.
\end{dinhly}

\begin{dinhnghia}
  \textbf{Tọa độ} của $v$ trong cơ sở $\mathcal{B}$: $[v]_{\mathcal{B}} = (x_1, \ldots, x_n)^T$ với $v = x_1 e_1 + \cdots + x_n e_n$.
\end{dinhnghia}

\begin{vidu}
  Cơ sở chính tắc của $\mathbb{R}^n$: $e_1 = (1,0,\ldots,0)^T$, $\ldots$, $e_n = (0,\ldots,0,1)^T$. $\dim \mathbb{R}^n = n$.
\end{vidu}

%% ================================================================
%%  CHƯƠNG 4: ÁNH XẠ TUYẾN TÍNH
%% ================================================================
\section{Ánh xạ tuyến tính}

\begin{dinhnghia}
  Ánh xạ $T: V \to W$ (giữa hai KGVT) gọi là \textbf{tuyến tính} (hay \textbf{đồng cấu}) nếu:
  \begin{itemize}
    \item $T(u + v) = T(u) + T(v)$ (bảo toàn phép cộng),
    \item $T(\lambda u) = \lambda T(u)$ (bảo toàn nhân vô hướng).
  \end{itemize}
  Tương đương: $T(\alpha u + \beta v) = \alpha T(u) + \beta T(v)$ với mọi $\alpha, \beta \in \mathbb{R}$, $u, v \in V$.
\end{dinhnghia}

\begin{tinhchat}
  $T(\mathbf{0}) = \mathbf{0}$, $T(-v) = -T(v)$. Ảnh của không gian con là không gian con; nghịch ảnh của không gian con là không gian con.
\end{tinhchat}

\begin{dinhnghia}
  \textbf{Hạt nhân} (kernel): $\ker T = \{v \in V \mid T(v) = \mathbf{0}\} \subset V$. \textbf{Ảnh} (image): $\mathrm{Im}\,T = T(V) = \{T(v) \mid v \in V\} \subset W$. Cả hai đều là không gian con.
\end{dinhnghia}

\begin{dinhly}[title={Định lý chiều}]
  $\dim V = \dim(\ker T) + \dim(\mathrm{Im}\,T)$. Số $\dim(\ker T)$ gọi là \textbf{số khuyết}, $\dim(\mathrm{Im}\,T)$ gọi là \textbf{hạng} của $T$.
\end{dinhly}

\begin{dinhnghia}
  \textbf{Ma trận của ánh xạ tuyến tính}: Cho $T: V \to W$, cơ sở $\mathcal{B}_V = \{e_1, \ldots, e_n\}$ của $V$, $\mathcal{B}_W = \{f_1, \ldots, f_m\}$ của $W$. Ma trận $A$ cỡ $m \times n$ thỏa $[T(v)]_{\mathcal{B}_W} = A [v]_{\mathcal{B}_V}$ có cột $j$ là $[T(e_j)]_{\mathcal{B}_W}$.
\end{dinhnghia}

\begin{dinhly}
  Nếu $A$ là ma trận của $T$ trong cơ sở $\mathcal{B}_V$, $\mathcal{B}_W$ thì $\mathrm{rank}(A) = \dim(\mathrm{Im}\,T)$, $\dim(\ker T) = n - \mathrm{rank}(A)$.
\end{dinhly}

\begin{dinhnghia}
  $T$ \textbf{đơn cấu} (đơn ánh) $\Leftrightarrow$ $\ker T = \{\mathbf{0}\}$. $T$ \textbf{toàn cấu} (toàn ánh) $\Leftrightarrow$ $\mathrm{Im}\,T = W$. $T$ \textbf{đẳng cấu} $\Leftrightarrow$ vừa đơn cấu vừa toàn cấu (khi đó $\dim V = \dim W$).
\end{dinhnghia}

%% ================================================================
%%  CHƯƠNG 5: TRỊ RIÊNG, VECTƠ RIÊNG, DẠNG TOÀN PHƯƠNG
%% ================================================================
\section{Trị riêng và vectơ riêng}

\begin{dinhnghia}
  Số $\lambda \in \mathbb{R}$ (hoặc $\mathbb{C}$) gọi là \textbf{trị riêng} của ma trận vuông $A$ nếu tồn tại vectơ $x \ne \mathbf{0}$ sao cho $Ax = \lambda x$. Vectơ $x$ gọi là \textbf{vectơ riêng} ứng với $\lambda$.
\end{dinhnghia}

\begin{congthuc}[title={Phương trình đặc trưng}]
  $\lambda$ là trị riêng của $A$ khi và chỉ khi $\det(A - \lambda I) = 0$. Đa thức $p_A(\lambda) = \det(A - \lambda I)$ gọi là \textbf{đa thức đặc trưng}. Các trị riêng là nghiệm của $p_A(\lambda) = 0$.
\end{congthuc}

\begin{dinhnghia}
  \textbf{Không gian riêng} ứng với $\lambda$: $E_\lambda = \ker(A - \lambda I) = \{x \mid Ax = \lambda x\}$. Số chiều $\dim E_\lambda$ gọi là \textbf{bội hình học} của $\lambda$; bội của $\lambda$ trong đa thức đặc trưng gọi là \textbf{bội đại số}.
\end{dinhnghia}

\subsection{Chéo hóa}

\begin{dinhnghia}
  Ma trận $A$ \textbf{chéo hóa được} nếu tồn tại ma trận khả nghịch $P$ sao cho $P^{-1}AP = D$ là ma trận chéo. Khi đó các cột của $P$ là các vectơ riêng, các phần tử trên đường chéo của $D$ là các trị riêng tương ứng.
\end{dinhnghia}

\begin{dinhly}[title={Điều kiện chéo hóa được}]
  Ma trận $A$ cấp $n$ chéo hóa được khi và chỉ khi $A$ có đủ $n$ vectơ riêng độc lập tuyến tính $\Leftrightarrow$ tổng các bội hình học của mọi trị riêng bằng $n$ $\Leftrightarrow$ $A$ có $n$ trị riêng (kể cả bội) và mỗi trị riêng có bội hình học = bội đại số.
\end{dinhly}

\begin{phuongphap}[title={Cách chéo hóa}]
  \begin{enumerate}
    \item Giải $\det(A - \lambda I) = 0$ tìm các trị riêng $\lambda_1, \ldots, \lambda_k$.
    \item Với mỗi $\lambda_i$, giải $(A - \lambda_i I)x = 0$ tìm cơ sở của $E_{\lambda_i}$.
    \item Nếu tổng số vectơ tìm được $= n$: đặt $P$ là ma trận có cột là các vectơ riêng (theo thứ tự trị riêng); $D = P^{-1}AP$ là ma trận chéo.
  \end{enumerate}
\end{phuongphap}

\subsection{Dạng toàn phương}

\begin{dinhnghia}
  \textbf{Dạng toàn phương} trên $\mathbb{R}^n$: $Q(x) = x^T A x = \sum_{i,j} a_{ij} x_i x_j$, với $A$ ma trận đối xứng ($A^T = A$). Có thể giả sử $A$ đối xứng vì $(x^T A x)^T = x^T A^T x = x^T A x$. Ma trận $A$ gọi là \textbf{ma trận của dạng toàn phương}.
\end{dinhnghia}

\begin{dinhly}[title={Dạng chính tắc}]
  Mọi dạng toàn phương đều đưa được về dạng chính tắc $Q = \lambda_1 y_1^2 + \cdots + \lambda_n y_n^2$ bằng phép đổi biến trực giao $x = Py$ (với $P$ ma trận trực giao: $P^T P = I$), trong đó $\lambda_i$ là trị riêng của $A$.
\end{dinhly}

\begin{dinhly}[title={Luật quán tính Sylvester}]
  Số hệ số dương và số hệ số âm trong dạng chính tắc không phụ thuộc cách đưa về chính tắc. Gọi $(p, q)$ là \textbf{chỉ số quán tính} với $p$ = số dương, $q$ = số âm.
\end{dinhly}

\begin{dinhnghia}[title={Xác định dấu}]
  Dạng toàn phương $Q(x)$:
  \begin{itemize}
    \item \textbf{Xác định dương}: $Q(x) > 0$ với mọi $x \ne \mathbf{0}$ $\Leftrightarrow$ mọi trị riêng $> 0$.
    \item \textbf{Xác định âm}: $Q(x) < 0$ với mọi $x \ne \mathbf{0}$ $\Leftrightarrow$ mọi trị riêng $< 0$.
    \item \textbf{Nửa xác định dương}: $Q(x) \ge 0$ và tồn tại $x \ne \mathbf{0}$ sao cho $Q(x) = 0$.
    \item \textbf{Không xác định}: $Q$ nhận cả giá trị dương và âm.
  \end{itemize}
\end{dinhnghia}

\begin{phuongphap}[title={Tiêu chuẩn Sylvester}]
  Ma trận đối xứng $A$ xác định dương khi và chỉ khi tất cả các định thức con chính (góc trên-trái) đều dương: $\Delta_1 > 0$, $\Delta_2 > 0$, $\ldots$, $\Delta_n > 0$.
\end{phuongphap}

\end{document}
